% TOP_PROBLEM: {"problemId": "problem.linearity-of-mapping-class-groups", "canonicalTitle": "Complex linearity of closed-surface mapping class groups by genus", "releaseRank": 281, "primaryDomain": "geometry_topology", "primaryDomainLabel": "Geometry & topology", "displayStatus": "open", "releaseStatus": "open"}
% !TeX program = lualatex
% ATTEMPT_STATUS: unresolved
% Research continuation by GPT_6_Astra_Ultra, 27 September 2026.
\documentclass[11pt]{article}
\usepackage[margin=25mm]{geometry}
\usepackage{fontspec}
\setmainfont{Latin Modern Roman}
\usepackage{amsmath,amssymb,mathtools}
\usepackage{unicode-math}
\setmathfont{Latin Modern Math}
\usepackage{xurl,hyperref}
\hypersetup{colorlinks=true,urlcolor=blue}
\setcounter{secnumdepth}{0}
\setlength{\parindent}{0pt}
\setlength{\parskip}{5pt}
\setlength{\emergencystretch}{3em}
\begin{document}
\section*{\texorpdfstring{Complex linearity of closed-surface mapping class groups by genus}{Complex linearity of closed-surface mapping class groups by genus}}\label{281}
\subsection{Definitions and mathematical statement}
Let $S_g$ be a closed connected oriented surface and
$\operatorname{Mod}(S_g)$ its group of orientation-preserving homeomorphisms modulo isotopy. Complex linearity means existence of a faithful finite-dimensional representation
\[
 \rho_g:\operatorname{Mod}(S_g)\hookrightarrow\mathrm{GL}_{n(g)}({\mathbb{C}}).
\]
The selected problem is to determine precisely the set of genera $g\ge3$ for which such a representation exists. Its dimension may depend on $g$, and neither unitarity nor a uniform dimension bound is imposed. Excluding representations below one fixed dimension, or resolving only an isolated genus, does not by itself classify all genera as required.
\par\smallskip\textit{Formulation and concept references:} \href{https://link.springer.com/article/10.1007/s00208-026-03530-5}{[S1]}.

\subsection{Short English statement}
For exactly which genera at least three can the mapping class group of a closed oriented surface be represented faithfully by finite-dimensional complex matrices?
\subsection{Research attempt: finite-index reduction and the dimension obstruction}
\textit{GPT-6 Astra Ultra, 27 September 2026. Status: UNRESOLVED.}

The quantifiers in this problem matter: the dimension may depend on the genus, but for each genus there must be one finite-dimensional representation whose kernel is trivial. A family of representations separating elements one at a time is not enough. I prove a finite-index reduction, test the natural homology representation, and give a finitely generated countermodel to the inference from residual finiteness to complex linearity.

\paragraph{The homology representation has a visible infinite kernel.}
The action on integral first homology gives
\[
 \rho_H:\operatorname{Mod}(S_g)\longrightarrow
 \operatorname{Sp}(H_1(S_g,\mathbb Z))\subseteq\operatorname{GL}_{2g}(\mathbb C).
\]
A Dehn twist about a simple closed curve $c$ acts by the transvection
\[
 v\longmapsto v+\langle v,[c]\rangle[c],                    \tag{1}
\]
up to the consistent sign convention for a positive twist. If $c$ separates the surface, then $[c]=0$ and the action is the identity.

Choose an essential separating curve with positive genus on each side; such a curve exists for every $g\ge3$. Its Dehn twist has infinite order. One way to see this is to choose a simple closed curve $a$ intersecting $c$ essentially. The standard annulus calculation gives
\[
 i(T_c^n(a),a)=|n|\,i(c,a)^2\qquad(n\in\mathbb Z),          \tag{2}
\]
where $i$ is geometric intersection number [E1]. Each crossing arc winds $n$ additional times in an annular neighborhood, giving the intersection count. In particular $T_c^n$ is nontrivial for $n\ne0$, although (1) sends all these powers to the identity.

This proves that $\rho_H$ is not the desired faithful representation. It says nothing by itself about representations which do not factor through homology. Tensor powers, direct sums, duals and algebraic constructions made solely from $\rho_H$ continue to kill its kernel; they cannot restore information lost at the first map. Adding a genuinely different representation might detect a particular separating twist, but faithfulness would still require detection of every nonidentity element of the Torelli group and beyond.

\paragraph{Linearity is equivalent to linearity of a finite-index subgroup.}
The forward implication is restriction. For the converse, let $H\le G$ have finite index $r$ and suppose $\rho:H\to\operatorname{GL}_d(\mathbb C)$ is faithful. Choose representatives $t_1=1,t_2,\ldots,t_r$ for the left cosets of $H$. For $g\in G$, write uniquely
\[
 gt_j=t_{\sigma_g(j)}h_{g,j},\qquad h_{g,j}\in H.
\]
On $W=\bigoplus_{j=1}^r\mathbb C^d$, define $\widetilde\rho(g)$ to send the $j$th summand into the $\sigma_g(j)$th summand by $\rho(h_{g,j})$. The identity
\[
 h_{gk,j}=h_{g,\sigma_k(j)}h_{k,j}
\]
proves multiplicativity, so this is a representation of dimension $rd$.

If $\widetilde\rho(g)=I$, the permutation of summands fixes the coset $H$, hence $g\in H$. On the first summand its action is $\rho(g)$, which forces $g=1$. Thus the induced representation is faithful. No normality assumption on $H$ was used.

For mapping class groups, this allows a reduction to any fixed congruence-level subgroup, for example the kernel of the action on $H_1(S_g,\mathbb Z/m\mathbb Z)$. Its index is finite because the target group is finite. If one constructs a faithful representation of that subgroup in finite dimension, the preceding formula constructs one for the full mapping class group. This is a useful reduction, but the reduced group is still infinite and the construction of its representation is the unproved step. Merely passing to a torsion-free subgroup does not supply the missing matrices.

\paragraph{What residual finiteness actually provides.}
Mapping class groups are residually finite by the classical theorem of Grossman [E2]. For any finite list $g_1,\ldots,g_s\ne1$, choose finite quotients $q_i:G\to Q_i$ with $q_i(g_i)\ne1$. The product quotient separates every element on the list. Its regular permutation representation gives a finite-dimensional complex representation doing the same.

This establishes
\[
 \text{for every finite list, some finite-dimensional representation separates it}.
                                                               \tag{3}
\]
The desired statement has the order of quantifiers reversed: there must be one representation separating \emph{all} nonidentity elements. Taking the direct sum over every finite quotient produces a generally infinite-dimensional representation. Neither residual finiteness nor (3) bounds the required dimension as the list grows.

The failure of this implication persists for finitely generated groups. Let
\[
 L=\Big(\bigoplus_{j\in\mathbb Z}\mathbb Z/2\mathbb Z\Big)\rtimes\mathbb Z,
                                                               \tag{4}
\]
where the generator $t$ shifts the lamp positions. The group is generated by $t$ and one lamp at position zero. It is residually finite: for any $m$ map positions modulo $m$ and reduce the shift modulo $m$, obtaining a finite quotient
\[
 L\longrightarrow (\mathbb Z/2\mathbb Z)^m\rtimes\mathbb Z/m\mathbb Z.
\]
For an element with nonzero shift $n$, choose $m$ not dividing $n$. For an element with zero shift and a nonempty finite lamp support, choose $m$ larger than the spread of that support; no two positions merge, so the image remains nonzero. Every nonidentity element is separated.

Yet $L$ cannot embed in $\operatorname{GL}_d(\mathbb C)$ for any finite $d$. Its lamp subgroup contains an infinite elementary abelian $2$-group. Commuting complex involutions are simultaneously diagonalizable: each is diagonalizable because its minimal polynomial divides $(X-1)(X+1)$, and its eigenspaces are preserved by every commuting operator; iterating the decomposition gives a common eigenbasis in finite dimension. Their image therefore lies in the diagonal subgroup $\{\pm1\}^d$, which has at most $2^d$ elements. An infinite lamp subgroup cannot inject into it.

This is not a subgroup assertion about mapping class groups. It is an explicit counterexample to a proposed general lemma, including the often added assumption of finite generation. Mapping class groups have additional structure, but that structure must be used in any successful upgrade of (3).

\paragraph{Why large commuting twist families are inconclusive.}
A different tempting obstruction comes from disjoint curves: their Dehn twists generate a free abelian group of rank as large as $3g-3$. The size of this rank alone does not rule out finite-dimensional complex linearity. In fact every finite-rank free abelian group embeds even in $\operatorname{GL}_1(\mathbb C)$:
\[
 (n_1,\ldots,n_r)\longmapsto\prod_{j=1}^r p_j^{n_j},        \tag{5}
\]
where the $p_j$ are distinct rational primes. Unique factorization makes (5) injective. Complex matrices need not have entries from a fixed finite set, and the target is not required to be a compact group.

The involution argument for (4) worked because all eigenvalues were constrained to a fixed finite set. The powers of disjoint twists have infinite order, so that argument does not apply to them. Confusing these two cases would give a false rank obstruction. Likewise a theorem excluding dimensions below a genus-dependent threshold leaves open every larger finite dimension and cannot classify the genera in the problem.

\paragraph{Source scope and first unresolved step.}
The 2026 source [S1] supplies criteria for certain automorphism-group extensions in terms of finite simple quotients of the underlying group. Its criterion is not an assertion that the hypotheses have been verified for every closed-surface mapping class group. One must also respect the distinction between an automorphism group, an outer automorphism group and a preimage of a subgroup: linearity passes to subgroups, but an arbitrary quotient of a linear group need not remain linear.

The local checks verify the transvection symplectic identity, finite lamp quotient operations, and the induced-representation multiplication on a small group example. They do not verify an infinite family of mapping-class relations. The proved output consists of the finite-index equivalence, an explicit failure of the homology route, and a countermodel to the residual-finiteness shortcut. The first missing step is either a faithful representation in a fixed finite dimension for a specified $g\ge3$, with every kernel element excluded, or an obstruction applying in every finite dimension. No classification of genera is obtained here; the requested problem remains unresolved.

\subsection{Sources}
\begin{itemize}
\item[S1]T. Koberda and M. Pengitore, Linearity criteria for automorphism groups of malabelian groups, Mathematische Annalen 395, 100 (2026), DOI 10.1007/s00208-026-03530-5; compare D. Margalit (2018), Question 1.1.
\url{https://link.springer.com/article/10.1007/s00208-026-03530-5}.
\textit{Formulation source.}
\item[E1]B. Farb and D. Margalit, \textit{A Primer on Mapping Class Groups}, Chapter 3, Dehn twists and intersection numbers. Used for (1)--(2), not for a linearity theorem.
\url{https://doi.org/10.1515/9781400839049}.
\item[E2]E. K. Grossman, \textit{On the residual finiteness of certain mapping class groups}, Journal of the London Mathematical Society (2) 9 (1974), 160--164.
\url{https://doi.org/10.1112/jlms/s2-9.1.160}.

\end{itemize}
\end{document}
